% §14.2 The BHF Construction
%
% Cross-reference resolution status:
%
% DEFINED here:
%   \label{sec:bhf-construction}                 -- this section
%   \label{def:bhf}                              -- the BHF definition
%   \label{alg:bhf-insert-query}                 -- the insert/query pseudocode
%
% ALREADY DEFINED (chapter 5):
%   \Cref{def:bloom-filter}                      -- chapter 5 §5.2 Bloom filter
%   \Cref{alg:bloom-insert-query}                -- chapter 5 §5.2 pseudocode
%
% ALREADY DEFINED (chapter 8):
%   \Cref{thm:entropy-floor}                     -- chapter 8 §8.3
%
% ALREADY DEFINED (chapter 14, earlier sections):
%   \Cref{sec:optimality-question}               -- §14.1
%
% FORWARD REFS within chapter 14 (resolve in later Plan 14 tasks):
%   \Cref{sec:space-analysis}                    -- §14.3 (Task 4)
%   \Cref{sec:fpr-analysis}                      -- §14.4 (Task 5)
%   \Cref{sec:construction-time}                 -- §14.5 (Task 6)
%
% §14.2 states the BHF construction in full. Two ingredients (perfect
% hash phi and fingerprint table F); insert algorithm; query algorithm
% with pseudocode; structural FNR-zero argument; FPR preview pointing
% to §14.4; worked example at n=1000, fpr=0.01.

\section{The BHF Construction}
\label{sec:bhf-construction}

\Cref{sec:optimality-question} sketched the answer in two ingredients
and two algorithms. This section gives the construction in full,
with a definition, pseudocode for the insert and query operations,
the structural argument for $\fnrate = 0$, a preview of the
false-positive analysis that \Cref{sec:fpr-analysis} completes, and
a worked example that places the BHF's bit budget side by side with
the standard Bloom filter at the same target error rate. The
ingredients are short. The argument that puts them together into an
FNR-zero approximate set is shorter.

\subsection*{The two ingredients}

Fix a finite universe $\Universe$ and a designated set $A \subseteq
\Universe$ with $|A| = n$. The construction depends on two hash
functions of different character.

\smallskip
\noindent\emph{(i) The perfect hash $\phi$.} The first ingredient is
a \emph{perfect hash function}
\[
  \phi \;:\; A \;\longrightarrow\; [n], \qquad [n] = \{1, 2, \ldots, n\},
\]
that is injective on $A$. Since $|A| = n$ and the codomain has size
$n$, injectivity makes $\phi$ a bijection between $A$ and the
addressing range. A perfect hash is not a hash function in the usual
sense of mapping a large universe uniformly into a small range; it
is a data-dependent function constructed at build time from $A$
itself, designed so that no two elements of $A$ collide. The cost of
building $\phi$ and the space it occupies are computed in
\Cref{sec:space-analysis} and \Cref{sec:construction-time}. For
the moment $\phi$ is a black-box bijection $A \to [n]$, extended in
some way (the specific extension does not affect FPR) to a function
$\phi : \Universe \to [n]$ on the entire universe.

The distinction between an ordinary hash function and a perfect
hash function deserves emphasis, since it is the key structural
departure from the Bloom filter design. An ordinary hash function
maps a large universe into a small range and collides at the rate
the range-to-universe ratio allows: by the birthday paradox, even
$\sqrt{n}$ distinct inputs into a range of size $n$ produce
collisions with high probability. The Bloom filter accommodates
collisions by spreading each insertion across $k$ independent
hashes, so that no single collision produces a false negative; the
$k$-way conjunction at query time absorbs the per-hash collision
noise. The BHF takes a different route. It commits to $A$ at build
time and constructs a hash function bespoke to $A$, guaranteed
collision-free on $A$ by design. The collision-free property is what
lets the BHF replace a $k$-way conjunction with a single fingerprint
comparison: there is no per-hash noise to absorb, because there are
no collisions. Constructing $\phi$ is therefore the workload that
the BHF moves from query time (where Bloom pays for it in extra hash
evaluations and bits) to build time (where it pays once and amortizes
over the structure's lifetime).

\smallskip
\noindent\emph{(ii) The fingerprint table $F$ and fingerprint hash
$g$.} The second ingredient is a \emph{fingerprint table}
\[
  F \;:\; [n] \;\longrightarrow\; \{0, 1\}^r,
\]
an array of $n$ slots, each holding an $r$-bit string. The
fingerprint width is
\begin{equation}
\label{eq:bhf-r}
  r \;=\; \lceil -\log_2 \fprate \rceil,
\end{equation}
the smallest integer for which $2^{-r} \le \fprate$. The table is
populated by a \emph{fingerprint hash}
\[
  g \;:\; \Universe \;\longrightarrow\; \{0, 1\}^r,
\]
a hash function that we assume is uniformly distributed on its
codomain. Concrete instances of $g$ are the low-order $r$ bits of
SHA-256$(x)$, MurmurHash3 truncated to $r$ bits, or any other family
known to behave uniformly on the inputs the application encounters.
The uniformity assumption is the only hash-function input the FPR
analysis of \Cref{sec:fpr-analysis} needs.

The fingerprint hash $g$ is an ordinary hash function, not a perfect
one. It is allowed to collide freely on $A$ and on $\Universe$, and
indeed the entire FPR analysis rests on $g$ producing uniform
collisions at rate $2^{-r}$. The two hashes therefore play
complementary roles: $\phi$ never collides on $A$ and serves as the
addressing function for the table; $g$ collides at a controlled rate
and serves as the fingerprint that the table uses to discriminate
$A$-elements from non-members. The width $r$ is the only place the
target $\fprate$ enters the construction; nothing else in the
algorithm depends on it. Doubling the target precision (from
$\fprate$ to $\fprate^2$) adds one bit to $r$ and doubles the table
size, recovering exactly the relation the entropy floor predicts.

\subsection*{Construction}

Given $A$, $\phi$, and $g$, the BHF table is built in one pass over
$A$:
\begin{equation}
\label{eq:bhf-insert}
  \text{for each } a \in A: \quad F[\phi(a)] \;\coloneqq\; g(a).
\end{equation}
The bijectivity of $\phi$ on $A$ means each slot of $F$ is written
exactly once. After the pass, the table $F$ records, for every
element of $A$, the fingerprint $g(a)$ stored at the slot $\phi(a)$
reserved for $a$. No slot remains unwritten because $\phi$ is a
bijection onto $[n]$; the table is fully populated.

\begin{definition}[Bernoulli Hash Function]
\label{def:bhf}
The \emph{Bernoulli Hash Function} (BHF) for a designated set $A$
with parameters $\phi$, $g$, and $r = \lceil -\log_2 \fprate \rceil$
is the data structure $(\phi, F)$ together with the membership query
operation:
\begin{itemize}
  \item \textsc{Query}$(x)$: return $\top$ if $F[\phi(x)] = g(x)$,
    return $\bot$ otherwise.
\end{itemize}
After the construction pass \eqref{eq:bhf-insert}, the predicate
$x \mapsto \textsc{Query}(x)$ defines the membership report of the
approximate set $\ASet{A}$.
\end{definition}

The definition is short because the construction is short. The data
structure stores the perfect hash representation and the fingerprint
table; the query reads one slot of the table and compares one
$r$-bit string to another. There is no flat bit array, no $k$-way
conjunction, no occupancy parameter to tune. Where the Bloom filter
spread information about $A$ thinly across the bits of a shared
array, the BHF concentrates it: each element of $A$ owns one slot,
and the slot stores the fingerprint that distinguishes it from
non-members.

The concentration of information is what closes the gap. In the
Bloom filter, the information about whether $x \in A$ is distributed
across $k$ bits of a shared array; each bit is a noisy partial
indicator and the conjunction of $k$ reads recovers the membership
verdict. The encoding's per-bit entropy is maximized at half
occupancy, and the resulting bit budget is $\log_2 e$ times the
floor. In the BHF, the information about whether $x \in A$ is
concentrated in one $r$-bit slot, and the slot's contents are
compared as a unit to the fingerprint $g(x)$ the query computes.
The encoding's per-bit entropy is no longer the right quantity to
optimize: the bit budget is set directly by the fingerprint width
$r$, which equals (up to integer rounding) the entropy floor's
$-\log_2 \fprate$. The structural change from a flat shared array
to a perfect-hash-addressed table is what makes the per-bit
optimization vanish, replaced by a per-fingerprint sizing that
matches the floor by construction.

\Cref{alg:bhf-insert-query} writes the construction and query in
pseudocode.

\begin{algorithm}[ht]
  \caption{BHF construction and query}
  \label{alg:bhf-insert-query}
  \SetKwInOut{Input}{input}
  \SetKwInOut{Output}{output}
  \SetKwProg{func}{procedure}{}{}
  \SetKwProg{fn}{function}{}{}

  \func{\textsc{Build}($A$)}{
    construct perfect hash $\phi : A \to [n]$\;
    allocate fingerprint table $F$ of size $n$, each slot $r$ bits\;
    \For{each $a \in A$}{
      $F[\phi(a)] \gets g(a)$\;
    }
    \Return $(\phi, F)$\;
  }
  \BlankLine
  \fn{\textsc{Query}($x$)}{
    $i \gets \phi(x)$\;
    \lIf{$F[i] = g(x)$}{\Return $\top$}
    \Return $\bot$\;
  }
\end{algorithm}

\textsc{Build} runs in one pass over $A$ plus the cost of
constructing $\phi$. \textsc{Query} runs in time dominated by the
two hash evaluations and the $r$-bit comparison: constant in $n$ and
linear in $r$, where $r$ is at most a few bytes in practical
deployments. The query algorithm does not consult any auxiliary
structure beyond $\phi$ and $F$. The pseudocode is the entire
runtime story.

Contrast with the corresponding Bloom filter pseudocode of
\Cref{alg:bloom-insert-query}: Bloom's \textsc{Insert} evaluates $k$
hash functions and writes $k$ bits; Bloom's \textsc{Query} evaluates
$k$ hash functions and reads $k$ bits, returning $\top$ only if all
$k$ reads return $1$. The BHF replaces $k$ hash evaluations with $2$
(the perfect-hash $\phi$ and the fingerprint hash $g$) and replaces
the $k$-way conjunction with a single $r$-bit equality test. Query
time is therefore lower than the Bloom filter's by a factor of
$k / 2$, which is approximately $(m/n) \ln 2 / 2 \approx 0.347
(m/n)$ at the Bloom optimum. The BHF is faster per query and
denser per bit; the price is that the construction is static (no
in-place insertions after build) and the perfect hash takes
$O(n)$ time to construct. \Cref{sec:construction-time} discusses
the trade-off in detail.

\subsection*{Why the false-negative rate is exactly zero}

The structural argument for $\fnrate = 0$ mirrors the Bloom filter
argument of \Cref{def:bloom-filter} but is even shorter. Suppose
$x \in A$. By the construction pass \eqref{eq:bhf-insert}, the
table was populated by setting $F[\phi(a)] = g(a)$ for every
$a \in A$, including $a = x$. Therefore $F[\phi(x)] = g(x)$ at the
end of the build phase. The fingerprint table is never modified
after construction (BHF is a static data structure), so $F[\phi(x)]
= g(x)$ remains true at every subsequent query time. The query
\textsc{Query}$(x)$ then reads slot $\phi(x)$, finds the fingerprint
$g(x)$ stored there, computes $g(x)$ from $x$, sees the two strings
match, and returns $\top$.

The argument depends on three facts: the construction wrote
$g(\cdot)$ into the slot $\phi(\cdot)$ at build time; the table is
not modified after construction; the perfect hash is a function,
so $\phi(x)$ at query time is the same slot $\phi(x)$ at build time.
The first is the definition of \textsc{Build}. The second is the
static-set assumption. The third is the function property of $\phi$.
For every $x \in A$ and every parameter setting,
\[
  \fnrate \;=\; \Prob{\textsc{Query}(x) = \bot \Given x \in A} \;=\; 0.
\]
The BHF is one-sided, on the same axis of the perfect filter line
as the Bloom filter. Its only failure mode is the false positive.

\subsection*{False-positive rate, preview}

Consider a query for $x \notin A$. The perfect hash $\phi$ is
defined on the universe but only meaningful as a bijection on $A$;
for $x \notin A$, the value $\phi(x)$ is some slot in $[n]$, namely
$\phi(a^*)$ for the unique $a^* \in A$ with $\phi(a^*) = \phi(x)$.
(The construction of $\phi$ may extend it to the universe by any
rule; for the FPR analysis below, what matters is that $\phi(x)$
lands in $[n]$.) The query reads $F[\phi(x)] = F[\phi(a^*)] =
g(a^*)$, and compares it to $g(x)$. The comparison succeeds iff
$g(a^*) = g(x)$.

By the uniformity assumption on $g$, the two values $g(a^*)$ and
$g(x)$ are independent uniform random strings in $\{0, 1\}^r$, and
the probability they agree is
\[
  \Prob{g(a^*) = g(x)} \;=\; 2^{-r}.
\]
With $r = \lceil -\log_2 \fprate \rceil$ from \eqref{eq:bhf-r}, this
gives $\fprate_{\text{BHF}} \le 2^{-r} \le \fprate$ as the designer
required. The bound is achieved by construction: choose $r$ to make
$2^{-r}$ no larger than the target, and the FPR is bounded
accordingly. \Cref{sec:fpr-analysis} develops the argument
carefully, including the role the perfect-hash extension to
$\Universe \setminus A$ plays in the analysis. For the purposes of
the present section the calculation
\[
  \fprate_{\text{BHF}} \;\le\; 2^{-r} \;=\; 2^{-\lceil -\log_2
  \fprate \rceil} \;\le\; \fprate
\]
is enough to compare against the floor and against the Bloom filter
in the worked example below.

\subsection*{A worked example}

Fix a designated set of size $n = 1000$ and a target false-positive
rate $\fprate = 0.01$. The fingerprint width is
\[
  r \;=\; \lceil -\log_2 0.01 \rceil \;=\; \lceil 6.644 \rceil \;=\; 7
  \text{ bits},
\]
since $2^{-7} = 1/128 \approx 0.0078 \le 0.01 \le 2^{-6} = 1/64
\approx 0.0156$. The fingerprint table occupies
\[
  n \cdot r \;=\; 1000 \cdot 7 \;=\; 7000 \text{ bits} \;=\; 875
  \text{ bytes}.
\]
Adding the perfect-hash representation, a modern construction
(quantified in \Cref{sec:space-analysis}) contributes on the order of
$2$ to $3$ extra bits per element at $n = 1000$, so the BHF total is
roughly $9 \cdot 1000 = 9000$ bits or $1125$ bytes. The leading term
is the fingerprint table; the perfect-hash term is the lower-order
correction.

The standard Bloom filter at the same target is sized by
\Cref{prop:bloom-overhead} of chapter~8 §8.4:
\[
  m \;=\; n \cdot \log_2(1/\fprate) \cdot \log_2 e
  \;=\; 1000 \cdot 6.644 \cdot 1.4427
  \;\approx\; 9586 \text{ bits},
\]
or about $1199$ bytes. The Bloom filter spends approximately
$9.59$ bits per element; the BHF spends $7$ bits per element on
fingerprints plus the sublinear perfect-hash overhead. At $n = 1000$
the two structures occupy similar total space, the perfect-hash
correction nearly cancelling the BHF's fingerprint savings. The
picture changes at scale. \Cref{sec:space-analysis} works out the
comparison at $n = 10^6$, where the perfect-hash overhead per
element shrinks well below one bit and the BHF runs at approximately
$7$ bits per element against the Bloom filter's $9.59$, recovering
the full $2.6$-bit advantage the entropy floor allows.

The same comparison at the tighter target $\fprate = 10^{-6}$ is
worth running quickly, because it shows the asymptotic gap that
\Cref{sec:space-analysis} formalizes. With $\fprate = 10^{-6}$, the
fingerprint width is $r = \lceil -\log_2 10^{-6} \rceil = \lceil
19.93 \rceil = 20$ bits, the floor is $19.93$ bits per element, and
the Bloom filter at the same target needs $19.93 \cdot 1.4427
\approx 28.76$ bits per element. The BHF at $n = 10^6$ pays
approximately $20$ bits per element on fingerprints plus the
sublinear correction, against the Bloom filter's $28.76$. The
absolute savings is roughly $8.8$ bits per element, or about $1.1$
megabytes saved on a million-element set. The relative savings is
about $30\%$ of the Bloom filter's budget, in line with the
$\log_2 e \approx 1.44$ ratio that \Cref{prop:bloom-overhead}
measured.

The entropy floor itself sits at
\[
  -\log_2 \fprate \;=\; -\log_2 0.01 \;=\; 6.644 \text{ bits per element}.
\]
The BHF's $r = 7$ bits per element exceeds the floor by $7 - 6.644
\approx 0.356$ bits, the integer-rounding gap that $\lceil \cdot
\rceil$ introduces. \Cref{sec:space-analysis} returns to the
rounding cost. The point of the worked example is to register the
order of magnitude: at this target rate, the BHF's per-element
budget is one integer bit above the floor, while the Bloom filter's
is approximately $2.95$ bits above the floor. The BHF closes about
$87\%$ of the gap that \Cref{prop:bloom-overhead} measured.
